\chapter{Introduction}

\section{Introduction}
Colleges, universities and polytechnic institutes run a large number of daily activities. They admit students, organise departments, programs, sections and courses, take attendance, hold examinations, publish results, collect fees, run hostels and libraries, manage staff, keep stores and assets, and publish notices. In many institutions these activities are handled by separate tools, spreadsheets or paper registers, so the same information is entered many times, records do not agree with each other, and students and guardians have to visit the office or phone staff to learn simple facts such as their attendance, result or fee due.

This project, the \textit{Next-Generation AI-Driven Smart Education Management System}, addresses this problem by bringing the academic and administrative work of an institution into one web-based Enterprise Resource Planning (ERP) platform. The platform is built for many institutions at once: each institution registers, receives its own isolated space, and manages its own users, roles, records and settings, while a platform-level Super Admin reviews registrations and manages subscriptions.

Alongside the regular management work, the system adds artificial intelligence where it helps people directly. A message and voice assistant answers questions from students, guardians and staff in Bangla or English. It reads the asker's own records for questions such as attendance, result, fees due or hostel meals, answers general questions from documents and answers approved by the institution, and forwards any question it cannot answer to a responsible person instead of guessing. A separate AI service offers machine learning models that estimate a student's academic risk and forecast GPA.

The work follows a service-based approach. The system is made of four applications: a Next.js frontend, a Core API that holds all business rules and is the single entry point for the frontend, an AI service for predictions, and a Voice service for language and speech processing. The Core API and the Voice service share one PostgreSQL database, and the services are packaged with Docker Compose.

The scope of this project covers multi-institution registration and subscription, user and role management, academic structure, student, guardian and teacher records, attendance, examinations and marks, finance, hostel, library, human resources, inventory, notices, a student and guardian portal, the message and voice assistant, and the prediction service. The project also has clear limitations. About half of the modules in the administration panel are interface screens that still show built-in sample data and are not yet connected to the backend. The prediction models are trained on synthetic data generated when the service starts, not on real institutional records. SMS, email and telephone delivery work only after an institution configures its providers, and the language model runs on the server's processor, which makes long answers slower than they would be on a graphics processor.

The expected outcome is a working platform in which an institution can run its daily operations from one place, every user sees only what their role allows, and students and guardians can get reliable answers at any time without waiting for office staff. The following chapters explain the background, the proposed system, its requirements, the methodology, the design, the implementation and the future direction of the work.

\section{Motivation}
The main motivation for this project is the gap between the amount of information an institution already holds and how difficult it is for people to use that information. Attendance, marks, fee invoices, hostel rent and meal menus are recorded every day, yet a student who wants to know whether a fee is due, or a guardian who wants to see a child's result, usually has to contact the office. Office staff, in turn, spend much of their time answering the same questions again and again.

A second motivation is fragmentation. When finance, hostel, library, human resources and academic records are kept in separate tools, the same student or staff member exists in several places and the records drift apart. A single system with one database removes this duplication and lets each module use the same student, teacher and guardian records.

A third motivation is security and accountability. Educational records are personal and sensitive. The system needs to make sure that a teacher sees only teaching data, an accountant sees only finance data, a guardian sees only the linked child's records, and one institution can never see another institution's data. It also needs to record who did what, block repeated password guessing, and force users to replace temporary passwords. These needs encouraged a design in which every request is checked against the user's current role and permissions, read from the database on every request, and every query is limited to the user's own institution.

A fourth motivation is the practical use of artificial intelligence. A general chatbot that invents answers is not acceptable for an educational institution. This project therefore designs an assistant that only answers from verified sources: the asker's own records, documents uploaded by the institution, and answers approved by an administrator. It refuses requests to change data, and when it is not sure, it sends the question to the right person and learns the approved answer for the future. The assistant runs on a language model hosted by the institution itself through Ollama, so questions and documents do not have to leave the institution's server. Voice support, with speech recognition and speech synthesis in Bangla and English, makes the same help available to people who prefer to speak rather than type.

Finally, the project is motivated by the need for a platform that can serve many institutions at once. Online registration, module selection, price quotation and subscription invoicing allow the same system to be offered to different institutions, each with its own data and settings.

\section{Objectives}
The aim of this project is to design and build a secure, multi-institution education management platform that unifies daily institutional work and supports it with practical artificial intelligence.
The main objectives of the project are as follows:
\begin{itemize}
    \item \textbf{Unified multi-institution management:} To provide one platform on which many institutions can register, choose the modules they need, receive a price quote and, after approval by the Super Admin, run their own isolated space. This is important because it removes separate tools and lets the same system serve many institutions. It is achieved by storing an institution identifier with every record, limiting every query to the signed-in user's institution, and providing platform screens for registration review, tenant status and subscription invoices. The expected outcome is that each institution manages its own data without any access to another institution's records.

    \item \textbf{Secure, role-based access:} To make sure that every user can reach only the screens and actions their role allows. This matters because educational records are sensitive. It is achieved through password hashing with bcrypt, JSON Web Token sessions, locking an account for 15 minutes after five wrong passwords, forcing a change of temporary passwords, built-in roles (Super Admin, Institution Admin, Teacher, Accountant, Hostel Warden, Storekeeper, Staff, Student, Guardian), custom roles with chosen permission keys, and an audit log of important actions. The expected outcome is a system in which a disabled account, a changed role or a suspended institution takes effect immediately on the next request.

    \item \textbf{Complete academic and operational records:} To manage the academic structure (departments, programs, sessions, sections and courses), students, guardians and teachers, attendance, examinations and marks, fee invoices and payments, hostel rooms, rent, meals and visitors, the library, human resources with leave and payroll, inventory and assets, and notices. This is important because these records are the basis of all reporting and of the assistant's answers. It is achieved through dedicated modules in the Core API backed by a single PostgreSQL database. The expected outcome is that each piece of information is entered once and used by every module that needs it.

    \item \textbf{Self-service for students and guardians:} To give students and guardians a portal where they can see their profile, academic results and attendance, fees and invoices, notices, library loans and hostel meals, and where students can book or cancel meals. This reduces the load on office staff and gives families timely information. It is achieved through portal endpoints that return only the records linked to the signed-in student or guardian. The expected outcome is that common questions are answered without visiting or calling the institution.

    \item \textbf{Trustworthy AI assistance and prediction:} To provide a message and voice assistant that answers in Bangla and English only from verified sources, forwards unanswered questions to responsible staff, and adds administrator-approved answers to its knowledge; and to provide models that estimate academic risk and forecast GPA. This is important because AI answers must be reliable in an educational setting. It is achieved with a locally hosted language model and embedding model through Ollama, document text extraction and search, speech-to-text with faster-whisper, text-to-speech with MMS voices, and Random Forest and Ridge regression models in the AI service. The expected outcome is an assistant that helps users at any time without inventing answers, and prediction results that are offered to staff as guidance.
\end{itemize}

The expected outcomes of the project are:
\begin{enumerate}
    \item A working web platform on which an institution can manage its academic, financial, hostel, library, human resources, inventory and notice work from one administration panel.

    \item Strict separation of data between institutions and between roles, with every request checked against the user's current permissions.

    \item A student and guardian portal and an assistant that together answer common questions about attendance, results, fees, meals and applications at any time.

    \item A knowledge base that grows with every question answered by staff and approved by an administrator, so that the assistant answers more questions over time.
\end{enumerate}

\section{Project Outline}
This project consists of eight structured chapters, each addressing a major component of the system's development. \textbf{Chapter~I} presents the background, motivation, objectives and an overview of the project structure. \textbf{Chapter~II} provides the literature review, discussing existing research on multi-tenant platforms, access control, student performance prediction, retrieval-based question answering and speech technology, and identifying the gap this project addresses. \textbf{Chapter~III} describes the proposed system, its architecture, user roles and major modules. \textbf{Chapter~IV} outlines the system requirements, with a focus on the hardware and software needed to develop and run the system. \textbf{Chapter~V} explains the methodology: the service-based development approach, request authorisation and data isolation, the machine learning models, the retrieval and ranking method of the message assistant, question escalation, speech processing and testing. \textbf{Chapter~VI} presents the major system design diagrams: workflow diagram, flowcharts, E--R diagrams, use-case diagram, activity diagram, sequence diagrams, DFD Level 0, DFD Level 1 and DFD Level 2. \textbf{Chapter~VII} discusses the system implementation, technology stack, deployment and testing. Finally, \textbf{Chapter~VIII} includes the conclusion, future work and business prospects.

\section{Summary}
This chapter introduced the \textit{Next-Generation AI-Driven Smart Education Management System}. It described the problem of scattered tools and repeated questions in educational institutions, and the motivation to solve it with one secure, multi-institution platform supported by trustworthy artificial intelligence. It stated five objectives: unified multi-institution management, secure role-based access, complete academic and operational records, self-service for students and guardians, and reliable AI assistance and prediction. It also set out the scope and the current limitations of the work and gave an overview of the remaining chapters, which together explain how the system was planned, designed and built.
